
\chapter{Regions}

\section{Region Structure}
\textit{Regions} are areas within the game world that have a consistent geography,
political identity, or other unifying factor that makes the locations found within
considerable as a consistent body that is distinct from other areas of the world.

\vspace{0.4em}

Regions have Towns, quest locations (see Locations in Chapter 10: Adventures), 
and traveling routes between them.

\vspace{0.4em}

Regions, Towns, and the Routes between and within them may be controlled or influenced by
factions (see Chapter 8: Loyalty: Factions \& Religion) of various spheres of influence, 
often with overlapping territories and conflict between them, in service of the plot.


\section{Towns \& Amenities}
For the purposes of these rules, \textit{Town} refers to inhabited settlements of all sizes, 
from small hamlets and isolated strongholds to sprawling metropolises.

\vspace{0.4em}

Towns may contain sub-locations that are roughly categorized as \textit{Civics}, 
\textit{Lodging}, \textit{Amenities}, and \textit{Shops}, although overlap may exist between them.

\vspace{0.4em}

Towns and their individual locations are generally controlled by at least one faction,
and access may be restricted individually or as a whole according to the loyalty a character has with that faction.


\subsection{Civics}
\textit{Civics} locations provide the necessary functions of civilization in a town, 
and are generally controlled by the town's dominant faction. 
These include places such as the mayor's manor or governor's palace, the town hall, a courthouse, a jail, 
and a citadel or barracks with an armory.

\vspace{0.4em}

These locations may also be used to divide the town into \textit{Districts}, 
allowing further faction-driven compartmentalization and control of the town,
with locations such as walls and gates, and guard posts or checkpoints.

\vspace{0.4em}

Civics locations, therefore, primarily serve a plot and roleplaying function, 
rather than providing any unique gameplay mechanic such as found at the other kinds of location.


\subsection{Shops}
\textit{Shops} are locations which provide equipment for sale or trade, 
allowing characters to acquire new gear in exchange for currency, 
rather than creating it from scratch.

\vspace{0.4em}

Shops may be either general or thematic, the latter of which often exist parallel to amenities, 
integrated together into a single location.


\subsection{Lodging}
\textit{Lodging} refers to any shelter where a character can find a long rest.
For example, lodging may be rented at an inn, charitably hosted by hospitable NPCs or factions, 
or purchased from the town to become character-owned properties.

\vspace{0.4em}

Characters must secure lodging in order to level-up in a town. 

\begin{DesignerNote}
This requirement enables the level-up mechanic to represent the process of comfortably recuperating, 
reflecting on the lessons learned from previous missions, then internalizing those lessons as new growth and development.
\end{DesignerNote}


\subsection{Amenities}
Certain passive abilities (i.e., Skills) may be triggered when you 
level-up in a town that provides a compatible \textit{Amenity}:

\begin{Table}{l X}
  \TableHeader{Amenity} & \TableHeader{Associated Skills} \\
  Workshop & Crafting, Jeweling \\
  Library & Studying, Writing \\
  Laboratory & Alchemy \\
  Arcaneum & Arcana, Enchanting \\
  Hospital & Healing injuries and conditions \\
  Temple & Prayer (see Religion in Chapter 8) \\
\end{Table}


\newpage


\section{Traveling Routes}
\textit{Traveling Routes} connect two locations (Towns, quest locations), 
or, within large or sprawling locations, two sub-locations.

\vspace{0.4em}

Routes have a length that is measured in the number of days it takes to travel that route.
For each day of travel, the traveling character, or party of characters, 
must roll on the \textit{Encounter Table} for that route. 

\vspace{0.4em}

Encounter Tables describe the interactions the party will have along the route.

\begin{Example}
  A particularly dangerous route may have bandits or monsters on the encounter table,
  whereas a relatively safe route may have traveling merchants, or other non-violent NPCs,
  such as kittens for Sierra.
\end{Example}




